\documentclass[11pt]{article}
\usepackage[margin=0.85in]{geometry}
\usepackage{booktabs}
\usepackage{enumitem}
\usepackage{parskip}
\setlength{\parskip}{0.4em plus 0.1em minus 0.1em}
\usepackage{titlesec}
\usepackage{hyperref}
\hypersetup{colorlinks=true, linkcolor=blue, urlcolor=blue}

\titlespacing*{\section}{0pt}{1.0ex plus 0.2ex minus 0.2ex}{0.4ex plus 0.1ex}
\titlespacing*{\subsection}{0pt}{0.8ex plus 0.2ex minus 0.2ex}{0.3ex plus 0.1ex}

\title{\textbf{Professor Cubazoid's 3D Tetris}\\[0.3em]
\large Executive Summary}
\author{STA 561 Final Project\\[0.3em]
ZRTMRH \quad Nicholas Popescu \quad Isabelle Dorage \quad Rally Lin}
\date{}

\begin{document}
\maketitle
\thispagestyle{empty}

\section*{The Problem}

Imagine you have a pile of oddly shaped 3D puzzle pieces---each made from three, four, or five small cubes glued together---and you need to figure out whether they can be assembled into a perfect larger cube with no gaps or overlaps. If they can, you need to show exactly how. If they cannot, you need to say so confidently.

This is a deceptively hard problem. A small 3$\times$3$\times$3 cube has only 27 cells, but a 10$\times$10$\times$10 cube has 1,000 cells and hundreds of pieces, each rotatable into up to 24 orientations. The number of possible arrangements grows explosively---brute force is completely impractical beyond trivially small instances. Our challenge was to build a solver handling cubes from 3$\times$3$\times$3 up to 25$\times$25$\times$25 reliably and efficiently.

\section*{Our Approach}

We built a three-layer solving system, where each layer addresses a different scale of the problem:

\textbf{Layer 1: Exact Solver.}
For small-to-medium puzzles (up to 9$\times$9$\times$9), we use a well-known combinatorial technique called ``exact cover,'' implemented with a data structure known as Dancing Links. This method systematically tries every possibility, guaranteeing a correct answer. It is fast for small instances but slows down dramatically as puzzles grow.

\textbf{Layer 2: Neural Network Heuristic.}
Inspired by the DeepCube project (which used deep learning to solve Rubik's cubes), we trained a 3D neural network to evaluate partial puzzle states. Given a half-assembled cube, the network predicts the probability that the remaining pieces can complete it. This learned intuition guides a ``beam search''---a focused exploration that follows the most promising assembly paths rather than trying everything. The network was trained on thousands of examples generated by the exact solver, then improved through a self-play process called Autodidactic Iteration, where the network learns from its own successes and failures.

\textbf{Layer 3: Divide and Conquer.}
For large cubes (10$\times$10$\times$10 and above), we decompose the problem into smaller sub-problems. The key insight is that a large cube can be split into smaller rectangular blocks along each axis---for example, an 11$\times$11$\times$11 cube can be divided into blocks of size 6$\times$6$\times$6, 6$\times$6$\times$5, 6$\times$5$\times$5, and so on. Each block is then solved independently using the exact solver. This decomposition makes previously intractable puzzles solvable in minutes.

A size-gated orchestrator automatically selects the right strategy for each puzzle size, combining all three layers into a seamless pipeline.

\section*{Results}

We conducted extensive testing across all grid sizes from 3 to 24, totaling over 1,800 individual test cases:

\begin{itemize}[nosep]
  \item \textbf{Perfect accuracy on solvable puzzles.} Every solvable test instance was solved correctly and independently verified (every cell covered exactly once, every piece used, shape matching confirmed). Across 757 positive test cases, the solve rate was 100\%.

  \item \textbf{Correct rejection of unsolvable puzzles.} Out of 323 fault tests (intentionally unsolvable), zero false solutions were returned.

  \item \textbf{Scalable performance.} Solve times range from under a second for small cubes to a few minutes for 18$\times$18$\times$18 (5,832 cells). The approach works reliably up to 24$\times$24$\times$24 (13,824 cells).
\end{itemize}

\begin{table}[h]
\centering
\begin{tabular}{@{}llll@{}}
\toprule
\textbf{Grid Size} & \textbf{Volume} & \textbf{Typical Solve Time} & \textbf{Method Used} \\
\midrule
3$\times$3$\times$3 & 27 & $<$ 0.1 seconds & Exact solver \\
5$\times$5$\times$5 & 125 & 1--5 seconds & Exact solver \\
10$\times$10$\times$10 & 1,000 & 10--30 seconds & Block decomposition \\
14$\times$14$\times$14 & 2,744 & 2--5 minutes & Block decomposition \\
18$\times$18$\times$18 & 5,832 & 3--5 minutes & Block decomposition \\
\bottomrule
\end{tabular}
\end{table}

\section*{The Role of AI in Development}

We used Claude extensively as a development collaborator---not as a black box, but as an accelerator for implementation and testing. AI assistance was particularly valuable for implementing the Dancing Links data structure, exploring decomposition strategies, designing comprehensive test suites, and diagnosing performance bottlenecks. Crucially, human insight remained essential for guiding AI-generated code: for instance, AI profiling revealed that the puzzle generator spent 94\% of its time on connectivity checks with a 97\% failure rate, but it was human observation that identified \emph{why}---cells in the ``middle'' of the structure cause disconnections---and proposed the fix (prioritizing peripheral cells), yielding a 7$\times$ speedup. The combination of human algorithmic insight and AI-accelerated iteration allowed us to build a far more sophisticated system than would have been feasible in the available time.

\section*{Future Work}

Three directions could extend this project significantly:

\begin{enumerate}[nosep]
  \item \textbf{Faster test generation.} Our current puzzle generator becomes slow for grids above 18$\times$18$\times$18, taking up to 10 minutes per instance. A faster constructive generator would unlock routine testing at even larger scales.

  \item \textbf{Scale-invariant neural networks.} Our current neural model is trained only on 3$\times$3$\times$3 instances. Training on diverse grid sizes with our scale-invariant architecture (which uses global average pooling instead of fixed-size layers) could produce a single model that provides useful guidance at any puzzle size.

  \item \textbf{Integration of learning and search.} Currently the neural network and block decomposition operate independently. A promising direction would be to use learned heuristics to guide the piece-to-block allocation during decomposition, potentially reducing the number of retries needed for large puzzles.
\end{enumerate}

\end{document}
